En la mitad derecha de un montaje sobre una guía recta, \(0\le x\le L\), la componente de fuerza medida sobre un eje global orientado hacia la derecha es
\[
F_+(x)=F_0\frac{x}{L}\left(1-\frac{x}{L}\right),
\qquad F_0,L>0.
\]
Todas las fuerzas son paralelas a la guía. La mitad izquierda se obtiene por reflexión geométrica: se reflejan tanto las posiciones como las flechas de fuerza. El eje global usado para asignar el signo de la lectura se mantiene fijo. El patrón completo se repite con longitud espacial \(2L\).

Un instrumento informa de esa componente con signo, \(F(x)\), y otro del módulo de la fuerza, \(M(x)\). Construya ambas extensiones al periodo completo y elija para cada registro una serie de senos, de cosenos o completa en las frecuencias \(n\pi/L\), justificándolo desde la regla física. Calcule la media espacial de cada registro y sólo el primer coeficiente no nulo del registro con signo, en la convención usual \(a_0/2+\sum_{n\ge1}(a_n\cos(n\pi x/L)+b_n\sin(n\pi x/L))\).
